% Generated draft: requires mathematical and source review.
\documentclass[11pt]{article}
\usepackage[margin=25mm]{geometry}
\usepackage{fontspec}
\setmainfont{Latin Modern Roman}
\usepackage{amsmath,amssymb,mathtools,unicode-math}
\setmathfont{Latin Modern Math}
\usepackage{enumitem,needspace,xurl,hyperref}
\setlength{\parindent}{0pt}
\setlength{\parskip}{5pt}
\setlength{\emergencystretch}{3em}
\newcommand{\N}{\mathbb{N}}
\newcommand{\Z}{\mathbb{Z}}
\newcommand{\Q}{\mathbb{Q}}
\newcommand{\R}{\mathbb{R}}
\newcommand{\C}{\mathbb{C}}
\newcommand{\F}{\mathbb{F}}
\newcommand{\A}{\mathbb{A}}
\newcommand{\PP}{\mathbb{P}}
\newcommand{\E}{\mathbb{E}}
\newcommand{\eps}{\varepsilon}
\newcommand{\dd}{\,\mathrm{d}}
\begin{document}
\section*{UnsolvedMath ID: 701502}
\textbf{Source problem code:} catalogue entry
\subsection{Definitions and mathematical statement}
For a group \(G\), define its derived series by \(G^{(0)}=G\) and \(G^{(n+1)}=[G^{(n)},G^{(n)}]\), where \([H,H]\) is the subgroup generated by the commutators of elements of \(H\). The group is solvable if \(G^{(d)}=\{1\}\) for some integer \(d\geq 0\); its derived length \(\operatorname{dl}(G)\) is the least such \(d\). A group is \(k\)-generated if it has a generating set of cardinality at most \(k\). An embedding is an injective group homomorphism. The abelianization is \(G_{\mathrm{ab}}=G/[G,G]\); it is torsion-free if its only element of finite order is the identity.

\textbf{Part (a), the remaining question in the supplied record.} Is it true that for every integer \(l\geq 2\) and every finitely generated solvable group \(G\) with \(\operatorname{dl}(G)=l\), there exist a solvable group \(H\) and an injective homomorphism \(\iota:G\hookrightarrow H\) such that
\[
H=\langle x,y\rangle\quad\text{for some }x,y\in H,
\qquad \operatorname{dl}(H)=l+1?
\]
Here the target length is read as exactly \(l+1\), following the wording of part (a).

\textbf{Part (b), retained as an archived question.} For every integer \(l\geq 2\), does there exist a positive integer \(k(l)\) such that every finitely generated solvable group \(G\) of derived length \(l\) embeds in a group \(H\) generated by at most \(k(l)\) elements and satisfying
\[
H^{(l+1)}=\{1\}?
\]
The bound \(k(l)\) must be independent of \(G\). In this formulation, the source phrase ``\((l+1)\)-solvable'' is interpreted as derived length at most \(l+1\). The record does not explicitly specify whether exact derived length was intended in that phrase; its archived result states exact derived length \(l+1\).

\textbf{Supplied qualifications and reported results.} The comment dated 2021 reports that every countable solvable group \(G\) of derived length \(l\) whose abelianization is torsion-free embeds in a group generated by at most two elements and having derived length \(l+1\). The archive reports that every finitely generated solvable group of derived length \(l\) embeds in a group generated by at most four elements and having derived length \(l+1\), thus supplying \(k(l)=4\) for part (b). Both reports are attributed to V. A. Roman'kov in the supplied record. They are recorded here as source-reported results, without independent verification. The supplied status is \texttt{partially\_solved}, with part (b) archived as solved and part (a) remaining in the main section.
\subsection{Short English statement}
The remaining question asks whether every finitely generated solvable group of derived length \(l\geq 2\) embeds in a two-generated solvable group of derived length \(l+1\). The supplied record reports a positive result when the abelianization is torsion-free, even for countable groups, and archives the version allowing four generators as solved. These are supplied source assessments requiring human review.
\subsection{Sources}
\begin{itemize}
\item Local UnsolvedMath registry, permanent ID \(701502\), exact source code \texttt{catalogue entry}, title \emph{Kourovka Notebook Problem 18.73}; supplied complete source record, upstream revision \texttt{d371ee917ac9647fbb06cd6fc529deec6f3dda82}.
\item A. Yu. Olshanskii, proposer of Problem 18.73 (2014), \emph{The Kourovka Notebook (Unsolved Problems in Group Theory)}, 18th issue (2014), updated in the 21st issue (2026), supplied identifier \texttt{arXiv:1401.0300v48}. The registry identifies the collection as Issues 1--21 (1965--2026), edited by E. I. Khukhro and V. D. Mazurov. The supplied material includes the main question, the comment of 2021, and the Archive of Solved Problems entry for part (b).
\item V. A. Roman'kov, \emph{Proc. Amer. Math. Soc.} \textbf{149} (2021), 4133--4143, as cited in the supplied comment and archive. The record supplies no article title, DOI, or article URL.
\item The supplied status review is dated \texttt{2026-10-16}, labels the problem \texttt{partially\_solved}, identifies the evidence as published, and places part (b) in the solved archive while retaining the remaining part in the main section. This is supplied dated curation, not an independently verified literature assessment.
\end{itemize}
\end{document}
